% ============================================================================
% CAPÍTULO 3: DESARROLLO DE LA PROPUESTA
% ============================================================================
% Este capítulo detalla el desarrollo práctico de la solución propuesta, la
% implementación del software, arquitectura del sistema, base de datos,
% o el núcleo del trabajo práctico o de ingeniería de tu tesis.
% ============================================================================

\chapter{Desarrollo de la Solución}

% --- Introducción al capítulo ---
Este capítulo detalla la arquitectura de software propuesta para el sistema de control de acceso vehicular de la Universidad Tecnológica Metropolitana, Sede Campus Ñuñoa. Dado el requerimiento de alta disponibilidad y la estricta restricción de baja latencia ($<100$ ms) para la lectura de patentes, la solución se diseña bajo un paradigma de computación distribuida. 

A continuación, se describen los dos nodos principales que componen el ecosistema: la plataforma de gestión centralizada (alojada en la nube) y el prototipo de inferencia local (ejecutado en el borde o \textit{Edge}), así como los mecanismos lógicos de sincronización de datos entre ambos. Cabe destacar que, al encontrarse el proyecto en la fase de propuesta, el diseño descrito a continuación representa la arquitectura teórica y el andamiaje lógico que guiará la implementación técnica del sistema en etapas posteriores de desarrollo.

\section{Aclaración de Alcance Arquitectónico}

Previo a detallar los componentes tecnológicos de la solución, es imperativo establecer una diferenciación conceptual entre el diseño teórico proyectado y el entregable material de esta investigación:
\begin{itemize}
    \item \textbf{Arquitectura de Referencia:} Corresponde al diseño a escala industrial (descrito en este capítulo) que contempla una red estructurada en la Nube y nodos físicos de hardware embebido (como NVIDIA Jetson) instalados en la portería del Campus.
    \item \textbf{Prototipo de Validación:} Corresponde a la implementación material acotada de este Trabajo de Título. Para validar la viabilidad lógica de la Arquitectura de Referencia y el rendimiento del modelo YOLOv12s, el sistema se construirá íntegramente en software y será simulado localmente en una computadora personal (PC), emulando el comportamiento del nodo de borde sin requerir instalación de hardware institucional en esta etapa.
\end{itemize}

\section{Arquitectura General (Topología Distribuida)}

Para aislar la carga de procesamiento computacional continuo (que exige la red neuronal) de las tareas administrativas cotidianas, la Arquitectura de Referencia divide el sistema físicamente en dos componentes interconectados:

\begin{itemize}
    \item \textbf{Nodo Central (Nube):} Encargado de la gestión administrativa institucional. Aloja el panel de control web principal y la base de datos maestra donde residen todos los registros históricos de acceso y usuarios.
    \item \textbf{Nodo de Borde (Edge / Prototipo Local):} Encargado exclusivamente de la captura de video en tiempo real, la inferencia de Inteligencia Artificial (YOLOv12s apoyado por EasyOCR) y la validación de ingreso. Está diseñado para operar localmente en la portería (simulado en PC para efectos de este prototipo) con el fin de garantizar tiempos de respuesta críticos sin depender de una conexión a internet constante.
\end{itemize}

\vspace{1cm}
\begin{figure}[H]
\centering
\begin{tikzpicture}[
    node distance=2.5cm and 2cm,
    box/.style={rectangle, draw=black, thick, fill=blue!10, text width=3.5cm, align=center, rounded corners, minimum height=1.5cm},
    user/.style={circle, draw=black, thick, fill=green!10, text width=2.2cm, align=center, minimum height=2.2cm},
    ext/.style={rectangle, draw=black, thick, fill=gray!20, text width=3.5cm, align=center, rounded corners, minimum height=1.5cm}
]

% Nodos
\node[user] (conserje) {Conserje / Admin};
\node[box, right=of conserje] (sistema) {\textbf{Sistema ALPR}\\ (UTEM Campus)};
\node[ext, right=of sistema] (porton) {Portón Vehicular\\ (Hardware)};
\node[user, below=of sistema] (vehiculo) {Vehículo / Funcionario};

% Flechas
\draw[->, thick, >=stealth] (conserje) -- node[anchor=south, font=\scriptsize] {Gestiona usuarios} (sistema);
\draw[->, thick, >=stealth] (vehiculo) -- node[anchor=east, font=\scriptsize] {Se aproxima (Video)} (sistema);
\draw[->, thick, >=stealth] (sistema) -- node[anchor=south, font=\scriptsize] {Abre/Cierra} (porton);

\end{tikzpicture}
\caption{Diagrama de Contexto (Nivel 1) del Sistema ALPR.}
\label{fig:c1_contexto}
\end{figure}

Al profundizar un nivel más en la arquitectura, la separación lógica entre ambos ambientes y sus contenedores de software internos se ilustra en la Figura \ref{fig:c2_contenedores}:

\vspace{1cm}
\begin{figure}[H]
\centering
\resizebox{\textwidth}{!}{
\begin{tikzpicture}[
    node distance=1.5cm and 2cm,
    layer/.style={rectangle, draw=blue, thick, fill=blue!5, text width=4cm, align=center, minimum height=1.2cm},
    db/.style={cylinder, draw=black, thick, shape border rotate=90, aspect=0.25, fill=gray!10, text width=2.5cm, align=center, minimum height=1.5cm}
]

% Nube (Central)
\node[layer] (spa) {\textbf{Panel Web}\\ (React SPA)};
\node[layer, below=of spa] (api) {\textbf{API Central}\\ (FastAPI)};
\node[db, right=of api, xshift=1cm] (pg) {\textbf{Base Maestra}\\ (PostgreSQL)};

\node[draw, dashed, thick, fit=(spa) (api) (pg), inner sep=0.5cm, label=above:\textbf{Nodo Central (Nube)}] (nube) {};

% Borde (Local)
\node[layer, below=of api, yshift=-1.5cm] (edge) {\textbf{Pipeline IA}\\ (Python + YOLO)};
\node[db, right=of edge, xshift=1cm] (sqlite) {\textbf{Caché Local}\\ (SQLite)};
\node[layer, left=of edge, xshift=-1cm] (cam) {\textbf{Cámara IP}\\ (Video Stream)};

\node[draw, dashed, thick, fit=(cam) (edge) (sqlite), inner sep=0.5cm, label=below:\textbf{Nodo de Borde (Portería)}] (borde) {};

% Flechas
\draw[->, thick, >=stealth] (spa) -- node[anchor=east, font=\scriptsize] {HTTP/REST} (api);
\draw[->, thick, >=stealth] (api) -- node[anchor=south, font=\scriptsize] {SQL} (pg);
\draw[->, thick, >=stealth] (cam) -- node[anchor=south, font=\scriptsize] {RTSP} (edge);
\draw[->, thick, >=stealth] (edge) -- node[anchor=south, font=\scriptsize] {Lee/Escribe} (sqlite);

% Sincronización
\draw[<->, thick, >=stealth, dashed] (api) -- node[anchor=east, font=\scriptsize] {Sincronización Asíncrona} (edge);

\end{tikzpicture}
}
\caption{Diagrama de Contenedores (Nivel 2) separando Nube y Borde.}
\label{fig:c2_contenedores}
\end{figure}



\section{Nodo Central (Plataforma Administrativa)}

El Nodo Central constituye el corazón administrativo del sistema. Está diseñado para ser desplegado en un servidor virtualizado (nube), garantizando disponibilidad ininterrumpida para el personal autorizado desde cualquier dispositivo con navegador web. Su arquitectura se compone de tres capas lógicas:

\subsection{Frontend: Interfaz Web Administrativa}
Se estructurará como una aplicación web de tipo SPA (\textit{Single Page Application}) construida con React y TypeScript. Esta interfaz proporcionará los módulos necesarios para que administradores y personal de conserjería puedan gestionar los perfiles de los funcionarios, matricular vehículos (asociando sus respectivas placas patentes) y visualizar en formato de tabla o gráfico los registros históricos de acceso al recinto.

\subsection{Backend: API RESTful Central}
La lógica de negocio transaccional será expuesta mediante una API REST desarrollada en Python con FastAPI. La elección de este framework se fundamenta en su capacidad nativa para manejar operaciones asíncronas de manera concurrente, lo que permite atender eficientemente múltiples peticiones simultáneas, tales como consultas al registro de ingresos o procesos de alta masiva de patentes.

\subsection{Almacenamiento: Base de Datos Maestra}
Se empleará PostgreSQL como motor de base de datos relacional principal. En esta base de datos se estructurará el modelo entidad-relación del proyecto, el cual mantendrá la integridad referencial entre la identidad de los funcionarios, la vigencia de sus vehículos y un \textit{log} inmutable de todos los accesos validados en el tiempo.

\section{Nodo de Borde (Edge)}

El Nodo de Borde representa la unidad computacional física que interactúa directamente con el entorno en la portería. El diseño arquitectónico recomienda para este fin hardware embebido con aceleración gráfica, como una NVIDIA Jetson Orin Nano. Sin embargo, para los fines de validación técnica de este estudio, dicho nodo será simulado temporalmente mediante un computador estándar (PC).

\subsection{Pipeline de Inferencia (Backend Local)}
El núcleo funcional de este nodo es un servicio local de alto rendimiento, orquestado mediante Python. Para cumplir con las estrictas cuotas de latencia exigidas, la ejecución del modelo de IA se optimiza utilizando las bibliotecas y \textit{bindings} de TensorRT (tal como se recomendó en el Marco Teórico), delegando la carga pesada directamente a la GPU. Este servicio ejecuta un ciclo de procesamiento (\textit{loop}) continuo para el análisis de video:
\begin{enumerate}
    \item \textbf{Captura:} Se obtiene el flujo de video en tiempo real de una cámara IP o circuito cerrado (CCTV) mediante la biblioteca OpenCV.
    \item \textbf{Preprocesamiento y Detección:} Se aplican filtros de mejora a la imagen y, posteriormente, la red convolucional YOLOv12s detecta y recorta espacialmente la placa patente dentro del encuadre.
    \item \textbf{Extracción OCR:} El recorte de la patente es procesado por EasyOCR para digitalizar los caracteres alfanuméricos.
    \item \textbf{Validación Lógica:} El texto digitalizado se contrasta inmediatamente contra la base de datos local para emitir un evento de autorización o denegación de acceso.
\end{enumerate}

\textbf{Selección del motor OCR para el Prototipo de Validación:} El Marco Teórico identifica a PaddleOCR como el motor de mayor rendimiento documentado, por lo que se mantiene como la alternativa recomendada para la Arquitectura de Referencia sobre NVIDIA Jetson y TensorRT. Para el Prototipo de Validación, en cambio, se adopta EasyOCR en virtud de tres consideraciones. En primer lugar, las limitaciones reportadas por la literatura (baja tasa de lecturas por minuto y confianza oscilante) corresponden a su ejecución en entornos carentes de aceleración gráfica, condición que no aplica al prototipo, el cual se ejecuta sobre un computador con GPU dedicada. En segundo lugar, su integración con detectores YOLO se encuentra documentada en sistemas ALPR de tiempo real que restringen la salida del OCR mediante verificación sintáctica de la patente; este proyecto extrapola dicho principio aplicando un filtro basado en los formatos de patente chilenos antes de consultar la base de datos local. En tercer lugar, EasyOCR opera sobre el mismo \textit{framework} de aprendizaje profundo (PyTorch) que el detector, lo que permite a ambos modelos compartir dependencias y memoria de GPU dentro de un único contenedor, reduciendo la complejidad de integración. Dado que el motor OCR se implementa como un módulo aislado dentro del \textit{pipeline}, esta elección será contrastada empíricamente con PaddleOCR durante el Trabajo de Título II, y podrá ser sustituida sin afectar al resto del sistema si no cumple las métricas de latencia y precisión de lectura definidas en el plan de validación.

\subsection{Visualización en Tiempo Real (Frontend Local)}
Para propósitos puramente de validación técnica y demostración del prototipo, el sistema de borde integrará una interfaz web ligera. Construida igualmente en React, pero servida de forma local, esta pantalla permitirá a los evaluadores visualizar en vivo el flujo de video procesado, remarcando con cuadros delimitadores (\textit{bounding boxes}) las detecciones de la red neuronal y demostrando la capacidad de inferencia del algoritmo.

\subsection{Persistencia Local (Caché \textit{Offline-First})}
Para lograr la latencia estricta de validación requerida y blindar el sistema ante eventuales caídas de la red institucional, el Nodo de Borde implementará una base de datos local utilizando SQLite. Esta base de datos actuará como un caché de alto rendimiento que almacenará el listado vigente de patentes autorizadas (para consulta rápida) y permitirá al propio sistema registrar de forma automática y local los accesos validados cuando opere sin conexión a internet, con el fin de sincronizarlos posteriormente hacia el nodo central una vez restablecido el enlace (sin requerir ni permitir intervención manual de los conserjes sobre este nodo).

\section{Mecanismo de Sincronización de Nodos}

Dado que la base de datos maestra (PostgreSQL) y la base de datos local de inferencia (SQLite) operan de forma desacoplada, el diseño contempla un mecanismo de sincronización asíncrono para mantener la coherencia de la información:

\begin{itemize}
    \item \textbf{Sincronización descendente:} Periódicamente, el servicio backend del nodo de borde consultará al nodo central en la nube para descargar e indexar en su SQLite local cualquier alta, baja o modificación ocurrida en la lista de vehículos autorizados.
    \item \textbf{Sincronización ascendente:} Cada vez que el nodo de borde certifique un acceso vehicular, el sistema almacenará automáticamente el registro (patente, fecha y hora) de forma local. En paralelo, y sin bloquear la barrera de acceso, este registro se enviará al nodo central para engrosar el histórico en PostgreSQL. Si la conexión a internet falla, el nodo local encolará automáticamente los registros en su memoria y los transmitirá en lote una vez restablecido el enlace, asegurando una arquitectura tolerante a fallos de forma transparente para el conserje.
\end{itemize}

\section{Modelo de Datos}

El modelo de datos se ha diseñado siguiendo principios de normalización relacional para garantizar la integridad referencial. A nivel macro, la base de datos maestra (PostgreSQL) centraliza la información completa, mientras que la base de datos local (SQLite) mantiene únicamente un subconjunto esencial de datos vigentes (las patentes autorizadas) para optimizar el rendimiento de las consultas en el borde.

\vspace{0.5cm}
\begin{figure}[H]
\centering
\begin{tikzpicture}[
    node distance=2cm and 2.5cm,
    entity/.style={rectangle, draw=black, thick, fill=blue!10, text width=4cm, align=center, rounded corners, minimum height=1.5cm},
    relation/.style={diamond, draw=black, thick, fill=yellow!20, text width=2cm, align=center, aspect=2}
]

% Entidades
\node[entity] (usuario) {\textbf{Funcionario}};
\node[entity, right=of usuario, xshift=2.5cm] (vehiculo) {\textbf{Vehículo}};
\node[entity, below=of vehiculo, yshift=-1.5cm] (acceso) {\textbf{Registro\_Acceso}};

% Relaciones intercaladas
\path (usuario) -- (vehiculo) node[midway, relation] (posee) {Posee};
\path (vehiculo) -- (acceso) node[midway, relation] (genera) {Genera};

% Conexiones
\draw[->, thick, >=stealth] (usuario) -- node[above, font=\scriptsize] {1} (posee);
\draw[->, thick, >=stealth] (posee) -- node[above, font=\scriptsize] {N} (vehiculo);
\draw[->, thick, >=stealth] (vehiculo) -- node[right, font=\scriptsize] {1} (genera);
\draw[->, thick, >=stealth] (genera) -- node[right, font=\scriptsize] {N} (acceso);

\end{tikzpicture}
\caption{Modelo Entidad-Relación conceptual del sistema ALPR.}
\label{fig:modelo_er}
\end{figure}

A nivel de diseño conceptual, las entidades principales son:
\begin{itemize}
    \item \textbf{Funcionario:} Representa al trabajador de la institución.
    \item \textbf{Vehículo:} Representa el automóvil autorizado, identificado inequívocamente por su placa patente, el cual pertenece a un funcionario.
    \item \textbf{Registro\_Acceso:} Tabla transaccional de tipo \textit{log} que guarda el historial exclusivo de los accesos exitosos (fecha y hora) para propósitos de auditoría.
\end{itemize}

\section{Casos de Uso y Flujos del Sistema}

Para comprender la operatividad del sistema integrado, se definen los tres flujos de trabajo principales que abordan el ciclo de vida completo de la información, desde el empadronamiento hasta el control físico en portería:

\begin{enumerate}
    \item \textbf{Gestión Administrativa (Alta de Vehículo):} El Administrador institucional ingresa al panel web del Nodo Central. Crea el perfil de un Funcionario e inscribe su vehículo asociado a su placa patente. Esta información se guarda en PostgreSQL. En el siguiente ciclo de sincronización, el Nodo de Borde descarga esta nueva patente y la indexa en su SQLite local, preparándose para leerla.
    \item \textbf{Inferencia Positiva (Acceso Autorizado):} Un vehículo llega al portón del Campus Ñuñoa. La cámara IP captura el video, el algoritmo YOLOv12s recorta la patente y EasyOCR la convierte a texto (ej. "AB-CD-12"). El script en Python consulta SQLite. Al encontrar la patente en estado activo, el sistema emite el pulso de apertura a la barrera y genera un registro temporal de acceso exitoso. Posteriormente, de forma asíncrona, este evento se envía al Nodo Central para alimentar el historial.
    \item \textbf{Inferencia Negativa (Acceso Denegado):} Un vehículo no registrado o transeúnte se aproxima al portón. El flujo de captura e inferencia se ejecuta normalmente, pero al consultar SQLite, la patente digitalizada no coincide con ningún registro activo. El sistema simplemente ignora la lectura y no emite señal a la barrera (permanece cerrada). En este escenario, \textbf{no se genera ningún registro en la base de datos} para evitar la saturación del almacenamiento (\textit{storage}) con vehículos ajenos a la institución que circulan por la calle o se detienen temporalmente frente al portón.
\end{enumerate}

\section{Diseño de Infraestructura Física y Hardware del Nodo de Borde}

\subsection{Infraestructura Preexistente y Supuestos de Motorización}

Al plantearse explícitamente como un proyecto de tipo \textit{Retrofit} (modernización inteligente de instalaciones existentes), el diseño del nodo de borde asume la disponibilidad y operatividad de la infraestructura pesada actual en la portería del Campus Ñuñoa. Esto excluye del alcance y del presupuesto del proyecto la adquisición de portones correderos, barreras vehiculares, motores de tracción y el tendido eléctrico de fuerza (220V). Cabe destacar que, para efectos de replicabilidad del diseño, cualquier portón o barrera física estándar que cuente con dimensiones adecuadas para el paso vehicular y un motor con capacidad de arrastre acorde a su peso, resultaría idóneo para la implementación. El sistema ALPR propuesto se acopla de manera universal y no invasiva a las entradas de control estándar (típicamente terminales de accionamiento manual o relé) de cualquier motorización actual, tal como se encuentra consignado en las limitaciones de alcance de la investigación.

\subsection{Justificación de Cómputo (NVIDIA Jetson)}

En el marco del diseño físico (\textit{Retrofit}) para el Nodo de Borde, la selección del componente de procesamiento central es crítica tanto para la viabilidad técnica como para la factibilidad económica del proyecto. Como se documentó en el Marco Teórico, la literatura avala el uso de arquitecturas ligeras (\textit{lightweight}) como la variante YOLOv12s (Small) para tareas de inferencia perimetral, la cual exige recursos de memoria altamente limitados (típicamente inferiores a 2 GB de VRAM para un tamaño de lote o \textit{batch} de 1).

Atendiendo a esta restricción paramétrica y a los lineamientos de optimización institucional, el diseño estipula el uso del \textbf{NVIDIA Jetson Orin Nano Developer Kit en su versión de 8GB} unificada. Esta elección se fundamenta en los siguientes pilares de diseño:

\begin{enumerate}
    \item \textbf{Holgura Operativa Segura:} Los 8 GB de RAM proporcionan un entorno holgado que absorbe sobradamente el consumo de la red neuronal (YOLOv12s) operando de forma concurrente sobre dos flujos de video independientes, el motor de reconocimiento óptico (EasyOCR), el sistema operativo base y los microservicios locales en Python, sin riesgo de agotamiento de memoria (\textit{Out Of Memory}). La literatura demuestra que configuraciones de 4 GB son suficientes para flujos individuales; por tanto, escalar a 8 GB garantiza una resiliencia extrema en producción ininterrumpida (24/7) para una arquitectura de doble vía.
    \item \textbf{Prevención de Sobredimensionamiento (\textit{Overprovisioning}):} Si bien existen en el mercado alternativas de grado industrial superiores, como la variante Jetson Orin NX de 16 GB, su implementación para este caso de uso de doble cámara representa un sobredimensionamiento técnico. Adquirir 16 GB de RAM para una red \textit{Small} implica inmovilizar un capital financiero en recursos de hardware que la aplicación \textit{nunca} llegará a utilizar.
    \item \textbf{Integridad del Flujo de Caja:} Al seleccionar el modelo de 8GB, se protege el Retorno de Inversión (ROI) y el Valor Actual Neto (VAN) del proyecto, garantizando que los fondos institucionales de la Universidad Tecnológica Metropolitana se inviertan con máxima eficiencia ingenieril y financiera, sin comprometer un milisegundo de rendimiento.
\end{enumerate}

\subsection{Justificación Óptica (Cámaras LPR)}

Para abarcar el control de acceso bidireccional (vía de entrada y vía de salida), el diseño estipula la instalación de dos cámaras \textbf{Dahua ITC413-PW4D-IZ1}, equipadas con un lente varifocal de 2.7 a 12 mm. Esta elección se contrapone deliberadamente a variantes de teleobjetivo (como el modelo Z3 de 8 a 32 mm), fundamentado en la topología física del acceso vehicular. 

Al tratarse de una portería universitaria donde el vehículo debe detenerse por completo frente a la barrera (velocidad de aproximación inferior a 20 km/h y distancias menores a 10 metros), el lente de 2.7-12 mm ofrece un ángulo de visión amplio (hasta 92$^\circ$) que permite encuadrar perfectamente el carril a corta distancia. Una cámara con lente de teleobjetivo, por el contrario, presentaría un punto ciego (\textit{blind spot}) significativo en proximidad, impidiendo capturar la patente de vehículos detenidos justo frente al portón. Por consiguiente, el uso de ópticas de distancia focal corta (como la variante IZ1 referenciada) resulta ser la alternativa topológicamente idónea y económicamente optimizada para un control de acceso de detención obligatoria.

\subsection{Interfaz de Potencia (Relé Optoacoplado)}

Para materializar la apertura física del portón desde la lógica de software, el diseño incorpora un módulo de relé optoacoplado de 1 canal. Este componente se alimenta con los 5V de la placa, pero es operado mediante señales de disparo de 3.3V lógicos, asegurando total compatibilidad con los pines GPIO del procesador sin riesgo de sobretensión. Al existir un único portón de acceso bidireccional, esta configuración de una vía (1 canal con capacidad de corte de 10 Amperes) es suficiente para emitir la señal de apertura. La selección de un componente con aislamiento óptico es un requerimiento de seguridad eléctrica ineludible: al utilizar un haz de luz interno para activar el interruptor mecánico, se aísla galvánicamente el costoso hardware de la Jetson Orin de cualquier potencial pico de voltaje inductivo proveniente del motor del portón. Dado que el sistema actúa únicamente como un "contacto seco" (cerrando un circuito de señalización de bajísimo amperaje hacia la placa controladora), la capacidad de 10A es el estándar óptimo que garantiza fiabilidad sin incurrir en costos injustificados de relés industriales de alta potencia. Físicamente, la conexión de baja tensión entre los bornes del relé (ubicado al interior del gabinete) y la placa controladora del motor se realizará mediante un tendido de \textbf{cable de control paralelo (2x18 AWG)}, protegido por la canalización conduit disponible.

Es imperativo destacar que el diseño eléctrico estipula conectar los contactos de este relé en configuración paralela respecto al sistema de apertura manual preexistente (receptor de radiofrecuencia). Esta topología paralela asegura que el portón conserve, en todo momento, su funcionalidad de apertura por control remoto convencional. De esta forma, el sistema ALPR actúa como un accionador complementario, proveyendo un mecanismo de \textit{fallback} o contingencia manual indispensable para dar acceso a motocicletas, vehículos de emergencia o casos especiales excluidos deliberadamente del alcance del modelo de inteligencia artificial.

\subsection{Alojamiento y Protección Ambiental (Gabinete Metálico)}

Para asegurar la viabilidad de la instalación en exteriores y resguardar el hardware computacional, se contempla el uso de un gabinete estanco metálico con certificación IP65. Estratégicamente, este gabinete no se adosará al pedestal de las cámaras, sino que se anclará firmemente al muro perimetral (muro interior del estacionamiento). Esta ubicación aleja el "cerebro" del sistema del portón, mitigando riesgos de sabotaje y facilitando enormemente la canalización de energía eléctrica (220V) desde el interior del edificio. Las dimensiones geométricas seleccionadas (250x200x150 mm) no son arbitrarias, sino que garantizan el holgado acomodo de la NVIDIA Jetson, el Switch PoE y el relé de control sobre la placa de montaje interior. Además, la profundidad de 150 mm provee el espacio necesario para alojar los transformadores de corriente y tomacorrientes sin estrangular el cableado, garantizando el volumen de aire necesario para la correcta disipación térmica pasiva del equipo de cómputo frente a variaciones meteorológicas y protegiéndolo íntegramente de polvo y humedad.

\subsection{Energización y Red (Switch PoE y Cableado UTP)}

Dado que las dos cámaras Dahua ITC413 especificadas para este diseño operan bajo el estándar PoE (\textit{Power over Ethernet}), se hace indispensable incorporar un Switch administrable PoE Gigabit (modelo DH-CS4006-4GT-60) y cableado UTP Cat6 para exterior en el presupuesto del nodo de borde. Cada cámara LPR, al activar sus iluminadores nocturnos a máxima potencia, exige un consumo pico cercano a los 20W. La selección de este switch en particular se fundamenta en su capacidad de suministro (\textit{Power Budget} de 60W), lo cual otorga un margen de seguridad holgado frente al consumo combinado de 40W de la óptica, previniendo apagones por sobrecarga. Asimismo, sus puertos Gigabit (1000 Mbps) aseguran que los flujos de video de alta resolución se transmitan hacia la NVIDIA Jetson sin cuellos de botella en la red local.

Respecto a la longitud del tendido de red, la topología física exige una canalización en tubería conduit de PVC subterránea desde el muro perimetral hasta el pedestal metálico. Según el levantamiento topográfico satelital (ver Figura \ref{fig:distancia_gabinete}), al reubicar el pedestal hacia el costado derecho del portón para evitar cualquier interrupción con la disposición del estacionamiento y el acceso vehicular, la trayectoria proyectada bordea los 9,3 metros. A modo de resguardo, considerando además la subida del cableado por el interior del poste de 3 metros, el diseño contempla la adquisición de 30 metros totales de cable UTP Cat6 100\% cobre para exterior (15 metros independientes por cada cámara). La inclusión de esta infraestructura pasiva y el switch permite transmitir simultáneamente los datos y la energía hacia las cámaras, simplificando la topología y centralizando la alimentación en la red eléctrica de 220V del gabinete, evitando instalar tomacorrientes expuestos en la calle.

\begin{figure}[H]
    \centering
    \includegraphics[width=0.6\textwidth]{figuras/distancia_gabinete.png}
    \caption{Levantamiento espacial de la distancia entre el muro perimetral (ubicación del gabinete) y el área de instalación del pedestal (9,27 m).}
    \label{fig:distancia_gabinete}
\end{figure}

\subsection{Estructura Metálica y Posicionamiento Óptico}

Respecto al anclaje y montaje físico, el levantamiento en terreno del pórtico (Campus Ñuñoa) evidenció que la infraestructura preexistente consta de una reja perimetral de barrotes metálicos y un área interior de estacionamiento con superficie de gravilla. Fijar equipamiento óptico directamente sobre esta reja es inviable por la oclusión visual ineludible de los barrotes, la vibración estructural del portón corredero, y el altísimo riesgo de vandalismo hacia el exterior de la calle. Por consiguiente, el diseño estipula la instalación de un único \textbf{pedestal metálico cónico de 3 metros de altura} posicionado estratégicamente por \textit{dentro} del campus. 

Aunque el poste posee una altura estructural de 3 metros, las cámaras LPR se instalarán mediante abrazaderas a una altura media (aprox. 1.2 a 1.5 metros) para garantizar un ángulo de captura óptimo hacia las patentes. Para habilitar una línea de visión despejada a esta altura, se contempla una intervención física consistente en el corte de un recuadro de 20x20 cm en los barrotes de la reja, a modo de ``ventana óptica''.

\subsection{Obra Civil y Canalización Subterránea}

Para asegurar la estabilidad física de esta estructura, el diseño de la obra civil incluye el despeje de la gravilla y el fraguado de un dado de fundación utilizando \textbf{hormigón preparado estructural de grado H20}, el cual fijará firmemente la placa base del pedestal mediante pernos expansivos. Para resolver la conectividad de datos, energía y señalización desde el gabinete perimetral, se tenderá una red de tuberías conduit de PVC subterráneas. Esta canalización contempla dos trayectos estratégicos: la ruta principal hacia el pedestal de las cámaras (con una extensión de 9.27 m) y una derivación que cruza el ancho del acceso vehicular hacia el motor del portón (ubicado en el flanco opuesto a la estructura del gabinete, como se documenta en el \textbf{Anexo~\ref{anexo:fotografias_terreno}}), protegiendo tanto los cables de red UTP Cat6 como el tendido de control del relé. La integridad de esta canalización se garantiza sellando todas sus uniones con adhesivo especial para PVC para prevenir filtraciones de humedad freática, empleando curvas de 90° y coplas que otorgan flexibilidad para sortear obstáculos (raíces o rocas) durante la excavación. 

La salida de los cables desde el gabinete metálico hacia el exterior se realiza por su cara inferior utilizando un terminal plástico de PVC con contratuerca, asegurando el tubo al chasis sin vulnerar su estanqueidad IP65. Finalmente, la tubería PVC viaja bajo la gravilla y queda ahogada en el centro del dado de hormigón, desembocando directamente en el interior hueco del poste. Esta topología unificada permite montar ambas cámaras (entrada y salida) espaldas con espaldas en un mismo pedestal blindado, aislando la óptica de vibraciones y ocultando el 100\% del cableado a la vista.
